\documentclass{article}
\usepackage{graphicx} % Required for inserting images

\title{Computer Science Lab Report}
\author{Eren Cicek - GH1053610}
\date{June 2026}
\begin{document}

\maketitle

\section{Project Links}
Portfolio Website:\\
https://erenci.github.io/computer-science-portfolio/\\
GitHub Repository:\\
https://github.com/erenci/computer-science-portfolio


\section{Introduction}
This report explains the design and creation of my portfolio website.
The website was created for the Computer Science Lab assessment. It presents information about my education, skills, and contact details. This report also describes the structure of the website and technologies used during the development process.
\section{Portfolio Website}
The portfolio website was created using HTML and CSS. It introduces my background and basic skills. The website contains six sections: About Me, Quick Facts, Skills, Education, Projects and Contact. I especially paid attention to using common standards in web development.
\section{GitHub Repository}
The GitHub repository contains the files used to create the project. It includes HTML and CSS files, a README file, my profile image and my LaTeX CV. I used Git to save my progress while developing the portfolio  and GitHub Pages to make my website available.  
\section{Design Decisions}
When designing the website, I wanted to keep everything simple. I chose a black background to make the text more readable. I also organised the content into six different sections. 
\section{Technologies Used}
I used HTML to build the structure of the website. Also CSS was used to improve its appearance. Git helped me keep track of changes during development, and GitHub was used to store the project online. GitHub Pages was used to publish the website, and I created my report and CV using LaTeX.
\section{Strengths and Weaknesses}
The only strength of my portfolio is its simple and organised structure. I think this project can be useful for beginners because it is easy to understand. However, the portfolio is still quite basic and does not include advanced features such as JavaScript or more advanced CSS styling. I also have only a small number of projects because I am still at the beginning of my studies.
\section{Improvement Ideas}
In the future, I would like to add more programming projects to my portfolio. I also plan to learn JavaScript and improve the overall design of the website. I would like to make the website more interactive and continue improving my web development skills.
\end{document}
